%============================================================
% CH1 · PATTERN 5 — EXERCISES, HINTS, ANSWERS
%============================================================
\subsection{Pattern 5: Exercises and Going Further}
\label{ssec:ch1-p5-ex}

\begin{patternbox}{Pattern 5 Summary Card}
\renewcommand{\arraystretch}{1.25}
\begin{tabularx}{\linewidth}{@{}>{\bfseries}l X@{}}
Recognise & one array \texttt{a[i]} = ``where $i$ goes next''; a deterministic process; ``eventually'', ``forever'', ``after $K$ steps'' \\
Shape & every walk is a $\rho$: tail $\mu$ + cycle $\lambda$, $\mu + \lambda \le n$; every piece has exactly one cycle (none if some \texttt{a[i] = -1}) \\
One start & walk with \texttt{first[v]} = step number; the first repeat gives $\mu$, $\lambda$ and the entry point; with \bigO{1} memory: Floyd's tortoise and hare (\hyperref[ssec:ch1-find-duplicate]{Problem~1.5.3}) \\
All vertices & three states (new / this walk / finished) or visit times $\ge$ start; or peel in-degree-0 vertices \\
Per vertex & cycle vertices first, then tails from the cycle outwards: $\text{ans}[v]$ from $\text{ans}[f(v)]$ \\
Top traps & fresh \texttt{visited} per start ($\bigO{n^2}$); old walk mistaken for a new cycle; \texttt{-1}; huge $K$ in \texttt{int}; deep recursion \\
\end{tabularx}
\end{patternbox}

\exlegend
\begin{enumerate}[leftmargin=*, label=\textbf{E5.\arabic*}, itemsep=6pt]
    \item \exeasy\ \pconcept \exprob{Array Nesting}{LeetCode 565}{https://leetcode.com/problems/array-nesting/}\ \lc{Medium}. \texttt{nums} is a permutation of $0..n-1$. The set $s[k] = \{\texttt{nums}[k], \texttt{nums}[\texttt{nums}[k]], \dots\}$ grows until an element would repeat. Return the largest $|s[k]|$. ($n \le 10^5$.)
    \item \exeasy\ \pconcept \exown{The $\rho$ of a generator.} The sequence $x_0 = 2$, $x_{i+1} = (x_i^2 + 1) \bmod 1000$ must eventually repeat. Find the tail length $\mu$ and the period $\lambda$ without storing more than one array of size $1000$.
    \item \exmedium\ \ptricky \exprob{Find Closest Node to Given Two Nodes}{LeetCode 2359}{https://leetcode.com/problems/find-closest-node-to-given-two-nodes/}\ \lc{Medium}. Out-degree $\le 1$ graph (\texttt{edges[i] = -1} allowed), two start vertices. Among the vertices reachable from both, return the one minimising the larger of the two distances (smallest index on ties), or $-1$.
    \item \exmedium\ \pclassic \exprob{Teleporter}{AtCoder ABC167 D}{https://atcoder.jp/contests/abc167/tasks/abc167_d}. Town $i$'s teleporter sends you to town $A_i$. Starting in town 1, where are you after exactly $K$ teleports? ($N \le 2\cdot10^5$, $K \le 10^{18}$.)
    \item \exmedium\ \pconcept \exown{How many shuffles?} A deck of $n$ cards is shuffled by always moving the card at position $i$ to position $p_i$ ($p$ is a permutation). After how many shuffles is the deck in its original order for the first time? ($p = [1,2,0,4,3] \to 6$.)
    \item \exhard\ \pclassic \cf{1700}\ \exprob{Mouse Hunt}{Codeforces 1027D}{https://codeforces.com/problemset/problem/1027/D}. A mouse in room $i$ moves to room $a_i$ every second; its starting room is unknown. A trap in room $i$ costs $c_i$. Find the minimum total cost that guarantees the mouse is eventually caught. ($n \le 2\cdot10^5$, $c_i \le 10^4$.)
    \item \exhard\ \ptricky \exprob{Maximum Employees to Be Invited to a Meeting}{LeetCode 2127}{https://leetcode.com/problems/maximum-employees-to-be-invited-to-a-meeting/}\ \lc{Hard}. Employee $i$ attends a round-table meeting only if seated next to \texttt{favorite[i]} ($\ne i$). Return the maximum number of employees that can be invited. ($n \le 10^5$.)
    \item \exhard\ \precent \cf{1700}\ \exprob{Medium Demon Problem (easy version)}{Codeforces 2044G1}{https://codeforces.com/problemset/problem/2044/G1}. Each of $n$ spiders starts with one plushie; every year each spider that has a plushie gives one to spider $r_i$ ($r_i \ne i$), all at the same time, and nobody keeps more than one. Year $y$ is stable if the distribution before year $y$ equals the one before year $y-1$. Find the first stable year. ($\sum n \le 2\cdot10^5$.)
\end{enumerate}

\subsubsection*{Hints}
\begin{enumerate}[leftmargin=*, label=\textbf{E5.\arabic*}]
    \item In a permutation every vertex also has in-degree 1, so there are no tails: $s[k]$ is exactly the cycle through $k$.
    \item The function $x \mapsto (x^2+1) \bmod 1000$ is a functional graph on $1000$ implicit vertices. Use Algorithm~1 (template \texttt{rho}).
    \item From each start, the reachable vertices form one path (no branching!). Record the distance along it, stopping at $-1$ or at a repeat.
    \item Store the walk from town 1 until the first repeat. Then $K$ is either inside the tail or can be reduced modulo the cycle length.
    \item Each cycle of $p$ returns to its start after exactly its length many shuffles, and only after multiples of it.
    \item Wherever the mouse starts, it ends up running around one cycle forever. Which rooms really need a trap?
    \item A seating is either one cycle of favourites on its own, or a collection of mutual pairs ($a \leftrightarrows b$), each extended by a chain of fans on both sides.
    \item Who loses the plushie first? A spider nobody gives to. Who next? Think of peeling the graph layer by layer.
\end{enumerate}

\subsubsection*{Answers}

\begin{solutionbox}[E5.1 Array Nesting]
Each walk from an unseen index is a complete cycle; a shared \texttt{seen} array makes the total work \bigO{n}.
\begin{lstlisting}[style=cpp]
#include <bits/stdc++.h>
using namespace std;

class Solution {
public:
    int arrayNesting(vector<int>& nums) {
        int n = nums.size(), best = 0;
        vector<bool> seen(n, false);
        for (int s = 0; s < n; s++) {
            if (seen[s]) continue;        // s already belongs to a measured cycle
            int len = 0, v = s;
            while (!seen[v]) {            // a permutation: the walk from s IS s's cycle
                seen[v] = true;
                v = nums[v];
                len++;
            }
            best = max(best, len);
        }
        return best;
    }
};
\end{lstlisting}
$[5,4,0,3,1,6,2]$: cycles $\{0,5,6,2\}$, $\{1,4\}$, $\{3\}$ $\Rightarrow$ \textbf{4}; \;$[0,1,2]$: three self-loops $\Rightarrow$ \textbf{1}.
\end{solutionbox}

\begin{solutionbox}[E5.2 The $\rho$ of a generator]
The sequence is $2, 5, 26, 677, 330, 901, 802, 205, 26, \dots$: the value $26$ first appears at step 2 and again at step 8, so $\mu = 2$ and $\lambda = 8 - 2 = 6$.
\begin{center}
\begin{tikzpicture}[x=1.25cm]
    \foreach \v/\i/\s in {2/0/gvsrc, 5/1/gv, 26/2/gvdst, 677/3/gvseen, 330/4/gvseen, 901/5/gvseen, 802/6/gvseen, 205/7/gvseen}{
        \node[\s, minimum size=9mm, font=\scriptsize\bfseries] (x\i) at (\i,0) {\v};
        \node[gdist, below=1mm of x\i] {step \i};
    }
    \foreach \i [evaluate=\i as \j using int(\i+1)] in {0,...,6} \draw[gde] (x\i) -- (x\j);
    \draw[gback] (x7.north) to[bend right=30] node[lbl, above]{$205^2 + 1 = 42026 \equiv 26$} (x2.north);
\end{tikzpicture}
\end{center}
\begin{lstlisting}[style=cpp]
#include <bits/stdc++.h>
using namespace std;
typedef long long ll;

int main() {
    const int m = 1000;
    vector<int> first(m, -1);             // implicit vertices 0..m-1, f(x) = (x*x + 1) mod m
    int x = 2, step = 0;
    while (first[x] == -1) {
        first[x] = step++;
        x = (int)(((ll)x * x + 1) % m);
    }
    cout << first[x] << " " << step - first[x] << "\n";   // 2 6
}
\end{lstlisting}
This is the ``$\rho$'' that gives Pollard's rho factorisation method its name. If even $m$ entries are too many, Floyd's tortoise and hare (template 5, used in \hyperref[ssec:ch1-find-duplicate]{Problem~1.5.3}) finds $\mu$ and $\lambda$ with \bigO{1} memory.
\end{solutionbox}

\begin{solutionbox}[E5.3 Find Closest Node to Given Two Nodes]
Out-degree $\le 1$ means the vertices reachable from a start lie on \emph{one} path, so ``BFS'' is just a walk. Compute $d_1$ and $d_2$ by walking (stop at $-1$ or a repeat --- the walk may enter a cycle), then scan $v = 0..n-1$ in increasing order.
\begin{center}
\begin{tikzpicture}
    \node[gvsrc] (a0) at (0,0.6) {0};
    \node[gvsrc] (a1) at (0,-0.6) {1};
    \node[gvdst] (a2) at (1.5,0) {2};
    \node[gv] (a3) at (3,0) {3};
    \draw[gde] (a0) -- (a2); \draw[gde] (a1) -- (a2); \draw[gde] (a2) -- (a3);
    \draw[-{Stealth[length=2mm]}, warningred, dashed] (a3) -- (4,0);
    \node[gdist, above=1mm of a2] {$d_1 = 1, d_2 = 1$};
    \node[gdist, above=1mm of a3] {$d_1 = 2, d_2 = 2$};
\end{tikzpicture}
\end{center}
\begin{lstlisting}[style=cpp]
#include <bits/stdc++.h>
using namespace std;

class Solution {
    vector<int> distFrom(const vector<int>& edges, int s) {
        vector<int> d(edges.size(), -1);  // -1 = not reachable from s
        for (int v = s, k = 0; v != -1 && d[v] == -1; v = edges[v], k++)
            d[v] = k;                     // stop at a dead end or when the walk repeats
        return d;
    }
public:
    int closestMeetingNode(vector<int>& edges, int node1, int node2) {
        vector<int> d1 = distFrom(edges, node1), d2 = distFrom(edges, node2);
        int best = -1, bestVal = INT_MAX;
        for (int v = 0; v < (int)edges.size(); v++) {  // increasing v: ties keep the smaller index
            if (d1[v] == -1 || d2[v] == -1) continue;
            int m = max(d1[v], d2[v]);
            if (m < bestVal) { bestVal = m; best = v; }
        }
        return best;
    }
};
\end{lstlisting}
$[2,2,3,-1]$, $0$, $1$ $\to$ \textbf{2} (distances 1 and 1); \;$[1,2,-1]$, $0$, $2$ $\to$ \textbf{2} (distances 2 and 0). Also checked against a brute force on $20\,000$ random graphs. Without the test \texttt{d[v] == -1} the walk would loop forever on a cycle.
\end{solutionbox}

\begin{solutionbox}[E5.4 Teleporter]
Record \texttt{order[t]} = town after $t$ teleports until the first repeat; with $\mu = \texttt{first}[v]$ and $\lambda = |\texttt{order}| - \mu$, the answer is \texttt{order[K]} if $K < \mu$, else \texttt{order}$[\mu + (K-\mu) \bmod \lambda]$.
\begin{lstlisting}[style=cpp]
#include <bits/stdc++.h>
using namespace std;
typedef long long ll;

int main() {
    int n;
    ll k;
    cin >> n >> k;
    vector<int> a(n + 1);
    for (int i = 1; i <= n; i++) cin >> a[i];

    vector<int> first(n + 1, -1);         // first[v] = step when v was first reached
    vector<int> order;                    // order[t] = town after t teleports
    int v = 1;
    while (first[v] == -1) {
        first[v] = order.size();
        order.push_back(v);
        v = a[v];
    }
    ll mu = first[v];                     // tail length (v = cycle entry)
    ll lam = (ll)order.size() - mu;       // cycle length
    if (k < mu) cout << order[k] << "\n";
    else        cout << order[mu + (k - mu) % lam] << "\n";
}
\end{lstlisting}
Sample 1 ($A = [3,2,4,1]$, $K = 5$): walk $1,3,4$ then $1$ again, so $\mu = 0$, $\lambda = 3$, answer \texttt{order}$[5 \bmod 3] = $ \textbf{4}. Sample 2 ($K \approx 7.3\cdot10^{17}$) prints \textbf{2}. Matches a step-by-step simulation on $5000$ random inputs with small $K$.
\end{solutionbox}

\begin{solutionbox}[E5.5 How many shuffles?]
A card on a cycle of length $L$ is back home after $t$ shuffles iff $L$ divides $t$. All cards are home iff every cycle length divides $t$: the answer is $\operatorname{lcm}$ of the cycle lengths.
\begin{lstlisting}[style=cpp]
#include <bits/stdc++.h>
using namespace std;
typedef long long ll;

ll shufflesToReturn(const vector<int>& p) {
    int n = p.size();
    vector<bool> seen(n, false);
    ll ans = 1;
    for (int s = 0; s < n; s++) {
        if (seen[s]) continue;
        ll len = 0;
        for (int v = s; !seen[v]; v = p[v]) { seen[v] = true; len++; }
        ans = ans / __gcd(ans, len) * len;   // lcm (may overflow for huge n: then work modulo primes)
    }
    return ans;
}

int main() {
    cout << shufflesToReturn({1, 2, 0, 4, 3}) << " " << shufflesToReturn({1, 2, 3, 0, 5, 6, 4}) << "\n";  // 6 12
}
\end{lstlisting}
$[1,2,0,4,3]$ has cycles of lengths 3 and 2 $\Rightarrow$ \textbf{6}; $[1,2,3,0,5,6,4]$ has 4 and 3 $\Rightarrow$ \textbf{12}. Note that the answer can be astronomically large for large $n$ --- real problems then ask for it modulo a prime, and the lcm is assembled from prime powers.
\end{solutionbox}

\begin{solutionbox}[E5.6 Mouse Hunt]
Wherever the mouse starts, it eventually runs around the cycle of its piece forever. A trap on a tail is useless for a mouse that starts on the cycle, while one trap anywhere on the cycle catches every mouse of that piece. So: \textbf{for each cycle, pay its cheapest room}; the answer is the sum. Use the three-state walk (theory, Algorithm 2) and walk once around each new cycle.
\begin{lstlisting}[style=cpp]
#include <bits/stdc++.h>
using namespace std;
typedef long long ll;

int main() {
    ios::sync_with_stdio(false);
    cin.tie(nullptr);
    int n;
    cin >> n;
    vector<int> c(n + 1), a(n + 1);
    for (int i = 1; i <= n; i++) cin >> c[i];
    for (int i = 1; i <= n; i++) cin >> a[i];

    vector<int> state(n + 1, 0);          // 0 new, 1 on the current walk, 2 finished
    ll total = 0;                         // up to 2e5 * 1e4 = 2e9: needs long long
    for (int s = 1; s <= n; s++) {
        if (state[s] != 0) continue;
        vector<int> path;
        int v = s;
        while (state[v] == 0) { state[v] = 1; path.push_back(v); v = a[v]; }
        if (state[v] == 1) {              // new cycle through v: trap its cheapest room
            int best = c[v];
            for (int u = a[v]; u != v; u = a[u]) best = min(best, c[u]);
            total += best;
        }
        for (int u : path) state[u] = 2;
    }
    cout << total << "\n";
}
\end{lstlisting}
The three official samples give \textbf{3}, \textbf{10} and \textbf{2}; the program also matches an exhaustive search over all trap sets on $5000$ random inputs with $n \le 8$. Note the \texttt{long long}: if every room points to itself, there are $2\cdot10^5$ cycles of cost up to $10^4$ each, a total of $2\cdot10^9 > 2^{31}-1$.
\end{solutionbox}

\begin{solutionbox}[E5.7 Maximum Employees to Be Invited to a Meeting]
The favourites form a functional graph. Two kinds of seatings work:
\begin{itemize}[leftmargin=*]
    \item \textbf{one cycle of length $\ge 3$}, seated in cycle order --- nobody else fits (each seat's neighbours are taken);
    \item \textbf{all 2-cycles at once}: for a mutual pair $a \leftrightarrows b$ seat the longest chain of fans ending at $a$, then $a, b$, then the longest chain ending at $b$; such blocks can be placed one after another around the table.
\end{itemize}
Answer $= \max(\text{longest cycle},\ \sum_{\text{2-cycles}} (\text{chain}[a] + \text{chain}[b]))$, where $\text{chain}[v]$ = vertices on the longest arrow chain ending at $v$ (including $v$). Peeling computes $\text{chain}$ because a vertex is always peeled before its successor.
\begin{center}
\begin{tikzpicture}
    \begin{scope}
        \foreach \v/\x/\y/\s in {0/0/0.7/gvsel, 3/1.4/0.7/gvsel, 4/1.4/-0.7/gvsel, 1/0/-0.7/gvsel, 2/-1.4/-0.7/gvdead}
            \node[\s] (e\v) at (\x,\y) {\v};
        \draw[gde] (e0) -- (e3); \draw[gde] (e3) -- (e4); \draw[gde] (e4) -- (e1); \draw[gde] (e1) -- (e0);
        \draw[gde] (e2) -- (e1);
        \node[note, text width=4.6cm] at (0,-1.7) {Example 3, $[3,0,1,4,1]$: one cycle $0 \to 3 \to 4 \to 1 \to 0$ of length 4 fills the table; $2$ cannot also sit next to $1$ $\Rightarrow$ \textbf{4}.};
    \end{scope}
    \begin{scope}[xshift=7.2cm]
        \foreach \v/\x/\y/\s in {4/-1.4/0.7/gvseen, 0/0/0.7/gvsel, 1/1.4/0.7/gvsel, 5/-1.4/-0.7/gvseen, 2/0/-0.7/gvsel, 3/1.4/-0.7/gvsel}
            \node[\s] (h\v) at (\x,\y) {\v};
        \draw[gde] (h4) -- (h0); \draw[gde] (h5) -- (h2);
        \draw[gde] (h0) to[bend left=30] (h1); \draw[gde] (h1) to[bend left=30] (h0);
        \draw[gde] (h2) to[bend left=30] (h3); \draw[gde] (h3) to[bend left=30] (h2);
        \node[note, text width=4.6cm] at (0,-1.7) {Own example $[1,0,3,2,0,2]$: two pairs, each with a fan chain of length 2 on one side: blocks $4,0,1$ and $5,2,3$ $\Rightarrow$ $3 + 3 =$ \textbf{6}.};
    \end{scope}
\end{tikzpicture}
\end{center}
\begin{lstlisting}[style=cpp]
#include <bits/stdc++.h>
using namespace std;

class Solution {
public:
    int maximumInvitations(vector<int>& favorite) {
        int n = favorite.size();
        vector<int> indeg(n, 0), chain(n, 1);  // chain[v]: longest chain ending at v
        for (int v = 0; v < n; v++) indeg[favorite[v]]++;
        queue<int> q;
        for (int v = 0; v < n; v++) if (indeg[v] == 0) q.push(v);
        while (!q.empty()) {                    // peel trees, push chains on
            int v = q.front(); q.pop();
            int w = favorite[v];
            chain[w] = max(chain[w], chain[v] + 1);
            if (--indeg[w] == 0) q.push(w);
        }
        int bigCycle = 0, pairs = 0;
        for (int s = 0; s < n; s++) {
            if (indeg[s] == 0) continue;        // not on a cycle, or already measured
            int len = 0, v = s;
            do { indeg[v] = 0; v = favorite[v]; len++; } while (v != s);
            if (len == 2) pairs += chain[s] + chain[favorite[s]];  // two mutual fans + their tails
            else bigCycle = max(bigCycle, len);
        }
        return max(bigCycle, pairs);
    }
};
\end{lstlisting}
The examples give \textbf{3}, \textbf{3}, \textbf{4}; an exhaustive search over all subsets and circular orders ($n \le 7$, $3000$ random inputs) agrees; the own example $[1,0,3,2,0,2]$ gives \textbf{6}. Why only \emph{one} long cycle but \emph{all} pairs? A cycle of length $\ge 3$ occupies both neighbours of each of its members, so the table is closed; a pair $a \leftrightarrows b$ is already happy with each other, leaving the outer seats free for fans --- and for the next block. This is the hardest exercise of the chapter; the peeling trick (``DP flowing along the trees towards the cycle'') returns in Chapter~11.
\end{solutionbox}

\begin{solutionbox}[E5.8 Medium Demon Problem (easy version)]
A spider has a plushie before year $y \ge 2$ iff someone who had one before year $y-1$ points to it. Cycle spiders always receive from their cycle predecessor, so they never lose theirs. A spider with in-degree 0 loses its plushie in year 1; a spider all of whose givers are gone after $k$ years loses it one year later. This is \textbf{peeling by layers} (theory, Algorithm 3): the spiders of layer $k$ hold a plushie before years $1..k$ and none before year $k+1$. So the distribution changes before every year up to $L+1$, where $L$ is the number of non-empty layers, and the first stable year is $\mathbf{L + 2}$. (Without tails, $L = 0$ and the answer is $2$: year 1 can never be stable.)
\begin{center}
\begin{tikzpicture}
    \foreach \v/\x/\y/\s/\t in {1/0/0.7/gvsel/{}, 2/0/-0.7/gvsel/{}, 4/1.6/-0.7/gvdead/{layer 3}, 3/3.2/-0.7/gvdead/{layer 2}, 5/4.8/-0.7/gvdead/{layer 1}}{
        \node[\s] (m\v) at (\x,\y) {\v};
        \node[gdist, below=1mm of m\v] {\t};
    }
    \draw[gde] (m1) to[bend left=30] (m2); \draw[gde] (m2) to[bend left=30] (m1);
    \draw[gde] (m5) -- (m3); \draw[gde] (m3) -- (m4); \draw[gde] (m4) -- (m2);
    \node[note, text width=4.8cm, align=left, anchor=west] at (5.6,0) {Sample test 3, $r = [2,1,4,2,3]$: three layers $\{5\}, \{3\}, \{4\}$, so the answer is $3 + 2 = 5$, as in the official explanation.};
\end{tikzpicture}
\end{center}
\begin{lstlisting}[style=cpp]
#include <bits/stdc++.h>
using namespace std;

int main() {
    int t;
    cin >> t;
    while (t--) {
        int n;
        cin >> n;
        vector<int> r(n + 1), indeg(n + 1, 0);
        for (int i = 1; i <= n; i++) { cin >> r[i]; indeg[r[i]]++; }
        vector<int> layer;                // current peeling layer
        for (int v = 1; v <= n; v++) if (indeg[v] == 0) layer.push_back(v);
        int layers = 0;
        while (!layer.empty()) {          // layer k loses its plushie in year k
            layers++;
            vector<int> next;
            for (int v : layer)
                if (--indeg[r[v]] == 0) next.push_back(r[v]);
            layer.swap(next);
        }
        cout << layers + 2 << "\n";       // last change in year layers + 1
    }
}
\end{lstlisting}
The official sample (five tests) prints \textbf{2, 2, 5, 4, 5}; a year-by-year simulation agrees on $20\,000$ random inputs with $n \le 9$. The \emph{tricky} part is the model, not the code: the statement talks about plushies and stability, but the answer is just ``number of peeling layers $+\,2$''.
\end{solutionbox}
